%==============================================================================
%  CAPÍTULO — LA INSOLVENCIA DE LOS GRUPOS EMPRESARIALES
%==============================================================================
\chapter{La insolvencia de los grupos empresariales}
\label{cap:grupos}

\fichaCapitulo{La laguna estructural del derecho concursal chileno.}{Ausencia de regulación de la
insolvencia de grupos, herramientas dispersas de la Ley \Nro 20.720, la técnica del levantamiento
del velo y las soluciones comparadas de consolidación procesal y sustancial.}

%==============================================================================
\bloqueTeorico
%==============================================================================

\subsection{El problema: una realidad económica sin estatuto propio}

La empresa contemporánea rara vez es una sola persona jurídica. La forma normal de organización de
la actividad económica de cierta escala es el \destacar{grupo empresarial}: una matriz que controla
filiales, coligadas y sociedades instrumentales, con patrimonios formalmente separados pero una
dirección económica unificada. Cuando ese grupo entra en crisis, la insolvencia no respeta las
fronteras societarias: la caída de la matriz arrastra a las filiales, los activos de una responden
por las deudas contraídas en beneficio de otra, y las garantías cruzadas entrelazan los pasivos.

El derecho concursal chileno afronta esa realidad con un instrumento pensado para el
\destacar{deudor individual}. La \LIR{} \citeyearpar{ley-20720} regula el concurso de \emph{una}
persona ---natural o jurídica---, no el de un conjunto de personas jurídicas vinculadas. No existe
en la ley un procedimiento de insolvencia de grupo, ni reglas de coordinación entre los concursos
de las distintas sociedades que lo integran. Es, probablemente, la \destacar{laguna estructural más
importante} del sistema, y la práctica de los grandes concursos la sufre a diario.

\begin{foco}[La asimetría entre la realidad y la norma]
El deudor se organiza como grupo para operar; el acreedor evalúa el riesgo del grupo para prestar;
el activo y el pasivo circulan dentro del grupo. Pero cuando llega la insolvencia, la ley obliga a
tratar cada sociedad como un compartimento estanco. Esa asimetría entre la unidad económica real y
la fragmentación jurídica del procedimiento es la fuente de casi todos los problemas que este
capítulo examina.
\end{foco}

\subsection{Por qué la separación patrimonial es a la vez virtud y problema}

La personalidad jurídica diferenciada y la consiguiente separación de patrimonios no son un defecto
que corregir: son la base del derecho societario y cumplen funciones valiosas ---permiten aislar
riesgos, atraer inversión y organizar actividades diversas---. El acreedor que contrató con una
filial contó con el patrimonio de \emph{esa} filial, y extenderle sin más el activo de las demás
sociedades del grupo defraudaría sus expectativas tanto como las de los acreedores de aquéllas.

El problema surge cuando la separación se \destacar{instrumentaliza}: cuando el grupo confunde
patrimonios en la operación pero los invoca separados en la crisis, cuando desplaza activos entre
sociedades para sustraerlos de determinados acreedores, o cuando fragmenta artificialmente una
empresa única en varias personas jurídicas para eludir reglas ---por ejemplo, la auditoría de la
reorganización ordinaria (capítulo \ref{cap:reorganizacion-simplificada})---. La cuestión, entonces,
no es abolir la separación patrimonial, sino \destacar{determinar cuándo el ordenamiento la
desconoce} en sede concursal.

%==============================================================================
\bloqueNormativo
%==============================================================================

\subsection{Las herramientas dispersas de la Ley \Nro 20.720}

A falta de un estatuto de grupo, el derecho chileno ofrece \destacar{piezas sueltas} que, bien
articuladas, permiten enfrentar parcialmente el problema:

\paragraph{El concepto de Persona Relacionada.} El \art{2} \Nro 26 define las personas relacionadas
por remisión al artículo 100 de la Ley \Nro 18.045 de Mercado de Valores, capturando a las
sociedades del grupo, sus controladores y administradores. Esa calificación produce efectos
concursales precisos: \destacar{posposición} de sus créditos (\art{63} y \art{241}),
\destacar{exclusión del quórum} de aprobación del acuerdo (\art{79}) y \destacar{plazo agravado} en
las acciones revocatorias (\art{287} y \art{288}). Es el principal instrumento con que la ley
reconoce la existencia del grupo, aunque sólo para \emph{neutralizar} su influencia, no para
consolidar su tratamiento.

\paragraph{Las acciones revocatorias.} Los actos entre sociedades del grupo celebrados en el período
sospechoso ---traspasos de activos, pagos preferentes, garantías cruzadas--- son atacables por la
vía de la ineficacia concursal (capítulo \ref{cap:revocatorias}), con el plazo ampliado que rige
para las personas relacionadas. Es la herramienta más eficaz para reconstruir el patrimonio drenado
por operaciones intragrupo.

\paragraph{La acumulación de procedimientos.} Aunque la ley no prevé un concurso único de grupo,
nada impide que los concursos de las distintas sociedades, radicados ante el mismo tribunal, se
coordinen de hecho, y que el mismo liquidador sea designado en varios de ellos ---lo que facilita la
gestión unificada sin fusionar las masas---.

\paragraph{La insolvencia transfronteriza.} Cuando el grupo es multinacional, el Capítulo VIII
permite coordinar los procedimientos abiertos en distintas jurisdicciones sobre las sociedades del
grupo (capítulo \ref{cap:transfronteriza}), aunque tampoco allí exista una consolidación
sustantiva.

\begin{alerta}[Lo que la ley reconoce y lo que ignora]
La Ley \Nro 20.720 \emph{reconoce} el grupo para restarle poder ---posponer y excluir a las
personas relacionadas--- y para revocar sus operaciones sospechosas. Pero \emph{ignora} el grupo
como unidad a administrar: no hay concurso consolidado, no hay reglas de reparto entre masas
vinculadas, no hay coordinación obligatoria de procedimientos. El reconocimiento es defensivo,
nunca organizativo.
\end{alerta}

\subsection{El levantamiento del velo en sede concursal}

Ante el vacío legal, la herramienta dogmática de cierre es el \destacar{levantamiento del velo}
(\emph{disregard of the legal entity}): la técnica que permite prescindir de la personalidad
jurídica formal cuando ésta se ha empleado de manera abusiva o fraudulenta, para alcanzar el
patrimonio que se pretende ocultar tras ella.

En Chile, el levantamiento del velo carece de consagración legal general y ha sido construido por la
jurisprudencia sobre bases dispersas ---la buena fe (\artcc{1546}), la prohibición del abuso del
derecho, la simulación, el fraude a la ley---. Sus presupuestos, exigentes, suelen resumirse en:

\begin{enumerate}
  \item \textbf{Control efectivo} de una sociedad sobre otra ---unidad de dirección real, más allá
  de la titularidad accionaria---.
  \item \textbf{Confusión de patrimonios} o de esferas de actuación ---cuentas comunes, activos
  usados indistintamente, ausencia de contabilidad separada real---.
  \item \textbf{Uso abusivo o fraudulento} de la separación ---la personalidad diferenciada se
  invoca precisamente para perjudicar a los acreedores o eludir obligaciones---.
  \item \textbf{Perjuicio} a terceros como consecuencia de ese uso.
\end{enumerate}

Concurriendo esos presupuestos, el juez puede \destacar{extender la responsabilidad} de una
sociedad a otra del grupo, o tratar como propios del deudor los bienes formalmente radicados en una
sociedad instrumental. En sede concursal, ello permite integrar a la masa activos que la fragmentación
societaria había sustraído.

\begin{practica}[El velo se levanta con prueba, no con sospecha]
El levantamiento del velo es una herramienta poderosa y, por lo mismo, de aplicación
\emph{excepcional} y estricta. Los tribunales lo conceden sólo ante prueba robusta de los cuatro
presupuestos, y rechazan su invocación como atajo para llegar a un patrimonio más solvente. Para el
liquidador o el acreedor que lo intenta, la clave es documentar la \destacar{confusión patrimonial}
concreta ---no la mera pertenencia al grupo--- y el \destacar{propósito fraudulento} de la
separación. Alegar «son todos lo mismo» sin acreditar la instrumentalización conduce al rechazo.
\end{practica}

%==============================================================================
\bloquePractico
%==============================================================================

\subsection{Las soluciones comparadas: consolidación procesal y sustancial}

Los ordenamientos que sí regulan la insolvencia de grupos distinguen dos técnicas de intensidad
creciente, y conviene manejarlas porque marcan el horizonte de una eventual reforma chilena.

\paragraph{La consolidación procesal (\emph{procedural consolidation}).} Es la forma leve: los
concursos de las distintas sociedades del grupo se \destacar{tramitan coordinadamente} ---mismo
tribunal, mismo administrador, calendarios y audiencias comunes, intercambio de información--- pero
\destacar{cada masa conserva su separación}. No se mezclan activos ni pasivos: sólo se gana
eficiencia en la administración. Es la solución que el \emph{Bankruptcy Code} estadounidense
practica de ordinario mediante la \emph{joint administration}, y la que la reforma brasileña de 2020
incorporó expresamente (capítulo \ref{cap:comparado-latam}).

\paragraph{La consolidación sustancial (\emph{substantive consolidation}).} Es la forma intensa y
excepcional: las masas de varias sociedades se \destacar{fusionan en una sola} ---un único activo
frente a un único pasivo, con extinción de los créditos y garantías intragrupo---. Sus efectos son
drásticos: el acreedor de una filial solvente pasa a concurrir con los de las filiales insolventes,
diluyéndose su expectativa de cobro. Por eso los ordenamientos que la admiten ---el estadounidense
por construcción jurisprudencial--- la reservan a casos de \destacar{confusión patrimonial extrema},
donde separar las masas es imposible o más costoso que fusionarlas, y donde ningún acreedor confió
razonablemente en el patrimonio aislado de una sola sociedad.

\begin{table}[htbp]
\centering
\small
\caption{Consolidación procesal y sustancial de grupos}
\label{tab:cap13c-consolidacion}
\begin{tabularx}{\textwidth}{@{}>{\raggedright\arraybackslash}p{0.22\textwidth}
                                 >{\raggedright\arraybackslash}p{0.34\textwidth} X@{}}
\toprule
\textbf{Rasgo} & \textbf{Consolidación procesal} & \textbf{Consolidación sustancial} \\
\midrule
Masas & Se conservan separadas & Se fusionan en una sola \\
\addlinespace
Créditos intragrupo & Subsisten & Se extinguen \\
\addlinespace
Efecto sobre el acreedor & Neutro: sólo cambia la administración & Puede diluir su cobro \\
\addlinespace
Presupuesto & Pertenencia al grupo y eficiencia & Confusión patrimonial extrema \\
\addlinespace
Frecuencia & Regla ordinaria donde se regula & Excepcional \\
\addlinespace
Estado en Chile & Posible de hecho (mismo liquidador y tribunal) & Sólo por vía del levantamiento
del velo \\
\bottomrule
\end{tabularx}
\end{table}

\subsection{Estrategia forense ante un grupo en crisis}

Mientras Chile no regule la materia, el abogado debe construir la solución con las piezas
disponibles. Las decisiones clave:

\begin{checklist}[Actuaciones ante la insolvencia de un grupo]
  \item \textbf{Mapear el grupo.} Estructura societaria, controlador, flujos de activos y pasivos
  intragrupo, garantías cruzadas. Es el insumo de toda decisión posterior.
  \item \textbf{Identificar las personas relacionadas} (\art{2} \Nro 26) para activar la
  posposición, la exclusión de quórum y el plazo revocatorio agravado.
  \item \textbf{Rastrear el período sospechoso} en busca de operaciones intragrupo revocables
  ---traspasos, pagos preferentes, garantías por deudas previas---.
  \item \textbf{Procurar la designación del mismo liquidador} en los concursos de las distintas
  sociedades, para una administración coordinada de hecho.
  \item \textbf{Evaluar el levantamiento del velo} sólo si hay prueba de confusión patrimonial y
  propósito fraudulento; reservarlo para la sociedad instrumental que concentra los activos
  desviados.
  \item \textbf{Coordinar la dimensión transfronteriza} (Capítulo VIII) si hay sociedades o activos
  del grupo en el extranjero.
\end{checklist}

\begin{foco}[La reforma pendiente]
La regulación de la insolvencia de grupos ---al menos en su forma procesal, coordinando los
concursos sin fusionar las masas--- es una de las asignaturas pendientes del derecho chileno, y de
las de menor costo político: no altera prelaciones ni expropia expectativas, sólo ordena la
administración de lo que hoy se administra de manera dispersa. La experiencia brasileña de 2020
ofrece un modelo cercano y probado. Mientras tanto, la coordinación depende de la iniciativa del
liquidador y de la benevolencia del tribunal, no de un derecho exigible.
\end{foco}
